%% abstract.tex: abstract in PT and EN  (FEUP regulations)
%% -------------------------------------------------------
\chapter*{Resumo}
%



\vspace{\fill}
\begin{flushleft}
{\Large \textbf{Palavras-chave}:} 
Ambiente de Borda \(\bullet\) 
Inferência Distribuída \(\bullet\)
Internet das Coisas \(\bullet\)
TinyML 
\end{flushleft}
\vspace*{24mm}

\acresetall %% to reset the acronym usage

\chapter*{Abstract}
% Here goes the abstract written in English. It should say the same.

In a world where \ac{AI} is increasingly integrated into everyday applications, the 
strategy for developing
and deploying state of the art \ac{AI} models has been relying heavily on powerful 
cloud-based servers. However, this approach has several drawbacks, including the
distance between the data source, the server and the end-user, which can lead to
increased latency, higher energy consumption, and potential privacy concerns.
To address these issues, edge computing has emerged as a promising solution, bringing
computational resources closer to where they are needed, in particular, the TinyML
paradigm, which provides a possibility to run \ac{AI} models on everyday devices.
Nonetheless, far-edge devices, such as microcontrollers, although
highly memory and computationally constrained, are widely used in the \ac{IoT} ecosystem,
equipped with sensors and actuators, and are capable of running small tasks while
consuming very low power. Faced with such devices, the possibility of running a
\ac{DNN} seems infeasible, but with the right strategies, it is possible to leverage
the capabilities of these devices and orchestrate them to achieve this task.
This work explores the potential of distributed inference,
where a \ac{DNN} model is partitioned and distributed across multiple
resource-constrained devices to achieve high performant inference in the edge 
environment. We present \ac{DisRUPT}, a distributed inference testbed, that enables the
evaluation of different partitioning strategies and deployment scenarios for \ac{DNN} 
models incorporating far-edge and edge devices together. Additionally, we present a set 
of strategies applied on different \ac{DNN} architectures, and evaluate their effects 
on inference latency and the network's energy consumption. From the collected results
in this work, we identify a set of best practices for distributed inference in far-edge
and edge devices, and, mainly, point out the existance of a input throughtput threshold 
where the inclusion of edge devices in the inference process becomes counterproductive.

\vspace{\fill}
\begin{flushleft}
{\Large \textbf{Keywords}:} 
Distributed Inference \(\bullet\)
Far-Edge Environment \(\bullet\) 
Internet of Things \(\bullet\)
TinyML 
\end{flushleft}
\vspace*{24mm}

\acresetall %% to reset the acronym usage
